%%%%%%%%%%%%%%%%%%%%%%%%%%%%%%%%%%%%%%%%%%%%%%%%%%%%%%%%%%%%%%%%%%%%%%
% Lớp 11 | Giới hạn. Hàm số liên tục | Giới hạn của hàm số
% Dạng 1D3?2-2: Thay số trực tiếp
% Số câu: 91  (sinh tự động bởi tach_cau.py; từ nay thầy sửa trực tiếp file này)
%%%%%%%%%%%%%%%%%%%%%%%%%%%%%%%%%%%%%%%%%%%%%%%%%%%%%%%%%%%%%%%%%%%%%%

%%%%%%%%%%%%%%%% TRẮC NGHIỆM 4 PHƯƠNG ÁN %%%%%%%%%%%%%%%%

% nguồn: Ok_Toan11_TracNghiem_Lan2
\begin{ex}%[1D3N2-2]%[ma:2978ba8f]
	$\lim\limits_{x\to 1}\left( x^2+3x+2 \right)$  có giá trị bằng
	 \choice
		 {$1$}
		{$2$} 
		{\True $6$} 
		{$+\infty $} 
	 \loigiai{
	Ta  có $\lim\limits_{x\to 1}\left( x^2+3x+2 \right)=1^2+3\cdot 1+2=6$.}
\end{ex}

% nguồn: Ok_Toan11_TracNghiem_Lan2
\begin{ex}%[1D3N2-2]%[ma:a124541b]
	Tính giới hạn $\lim\limits_{x\to  2}\dfrac{x+2}{x-1}$ ta được kết quả là
	\choice 
		{\True $4$} 
		{$1$} 
		{$2$} 
		{$3$} 
	 \loigiai{
	 Dễ thấy $\lim\limits_{x\to 2}\dfrac{x+2}{x-1}=\dfrac{2+2}{2-1}=4$.}
\end{ex}

% nguồn: Ok_Toan11_TracNghiem_Lan2
\begin{ex}%[1D3N2-2]%[ma:eba97e91]
	 Tính giới hạn $\lim\limits_{x\to 1}\dfrac{2x^3+3x-1}{x^2+1}$ ta được kết quả bằng
	 \choice
		{$1$} 
		{\True  $2$}
		{$3$} 
		{$4$} 
	\loigiai{ 
	Ta có  $\lim\limits_{x\to 1}\dfrac{2x^3+3x-1}{x^2+1}=\dfrac{{{2\cdot 1}^3}+3\cdot 1-1}{1^2+1}=\dfrac{4}{2}=2$.}
\end{ex}

% nguồn: Ok_Toan11_TracNghiem_Lan2
\begin{ex}%[1D3N2-2]%[ma:29875372]
	 Tính giới hạn $L=\lim\limits_{x\to 3}\dfrac{x-3}{x+3}$
	 \choice
		{$1$} 
		{$-\infty $} 
		{\True  $0$}
		 {$+\infty $}
	 \loigiai{
	Ta có  $L=\lim\limits_{x\to 3}\dfrac{x-3}{x+3}=\dfrac{0}{6}=0$.}
\end{ex}

% nguồn: Ok_Toan11_TracNghiem_Lan2
\begin{ex}%[1D3N2-2]%[ma:c6b973ae]
	Tính giới hạn $\lim\limits_{x\to  0}\left( 2x^2+3x-5 \right)$
	\choice 
		 {$0$}
		 {$3$}
		{$2$} 
		{\True  $-5$}
	\loigiai{ 
	Ta  có $\lim\limits_{x\to 0}\left( 2x^2+3x-5 \right)=-5$.}
\end{ex}

% nguồn: Ok_Toan11_TracNghiem_Lan2
\begin{ex}%[1D3N2-2]%[ma:2301bfa4]
	Tính  $\lim\limits_{x\to 1}\dfrac{x^2+3x+2}{-2x^2+x+3}$
	\choice 
		{$-\dfrac{1}{2}$} 
		 {$\dfrac{2}{3}$}
		{\True  $3$}
		{$\dfrac{1}{5}$} 
	 \loigiai{
	Ta có $\lim\limits_{x\to 1}\dfrac{x^2+3x+2}{-2x^2+x+3}=\lim\limits_{x\to 1}\dfrac{1^2+3\cdot  1+2}{-{{2\cdot 1}^2}+1+3}=3$.}
\end{ex}

% nguồn: Ok_TracNghiemLop11_Lan1
\begin{ex}%[1D3N2-2]%[ma:f376c0fe]
	Tính giới hạn $A=\lim\limits_{x \to 1} \dfrac{x^3-4x^2+2}{x^2-3x+3}$ 
	\choice 
		{\True  $-1$}
		{$-\infty$} 
		 {$\dfrac{4}{3}$}
		{$+\infty$} 
	 \loigiai{
		 Ta có $A=\lim\limits_{x \to 1} \dfrac{x^3-4x^2+2}{x^2-3x+3}=\dfrac{1-4+2}{1-3+3}=-1$.
	 }
\end{ex}

% nguồn: Ok_TracNghiemLop11_Lan1
\begin{ex}%[1D3N2-2]%[ma:80450750]
	Giá trị  của giới hạn $\lim\limits_{x\to 2}\sqrt[3]{\dfrac{x^2-x+1}{x^2+2x}}$ là
	 \choice
		{$\dfrac{1}{4}$} 
		 {\True $\dfrac{1}{2}$}
		 {$\dfrac{1}{3}$}
		{$\dfrac{1}{5}$} 
	 \loigiai{
		Ta có $\lim\limits_{x\to 2}\sqrt[3]{\dfrac{x^2-x+1}{x^2+2x}}=  \sqrt{\dfrac{2^2-2+1}{2^2+2 \cdot 2}}=\dfrac{1}{2}$.
	 }
\end{ex}

% nguồn: Ok_TracNghiemLop11_Lan1
\begin{ex}%[1D3N2-2]%[ma:d9a76c2d]
	Giá trị của giới hạn $\lim\limits_{x\to 3}\sqrt{\dfrac{9x^2-x}{(2x-1)(x^4-3)}}$ là 
	\choice 
		 {$\dfrac{1}{5}$}
		{$\sqrt 5$} 
		{\True  $\dfrac{1}{\sqrt 5}$}
		{$5$} 
	 \loigiai{
		Ta có $\lim\limits_{x\to  3}\sqrt{\dfrac{9x^2-x}{(2x-1)(x^4-3)}}
		 =\sqrt{\dfrac{9 \cdot 3^2-3}{(2 \cdot 3-1)(3^4-3)}}
		=\dfrac{1}{\sqrt 5}$. 
	 }
\end{ex}

% nguồn: Ok_TracNghiemLop11_Lan1
\begin{ex}%[1D3N2-2]%[ma:1050b828]
	Giá trị  của giới hạn $\lim\limits_{x\to\sqrt 3}|x^2-4|$ là
	\choice 
		{$0$} 
		{\True  $1$}
		 {$2$}
		{$3$} 
	\loigiai{ 
		Ta có  $\lim\limits_{x\to\sqrt 3}|x^2-4|=|(\sqrt 3)^2-4|=1$.
	} 
\end{ex}

% nguồn: Ok_TracNghiemLop11_Lan1
\begin{ex}%[1D3N2-2]%[ma:51a36838]
	Tính $\displaystyle\lim_{x\to-1}\left(2x^3-3x\right)$  bằng
	 \choice
		{$-1$} 
		 {$-5$}
		{\True $1$} 
		 {$5$}
	 \loigiai{
		$\lim\limits_{x\to-1  }\left(2x^3-3x\right)=2\cdot (-1)^3-3\cdot (-1)=1$.
	} 
\end{ex}

% nguồn: Ok_TracNghiemLop11_Lan1
\begin{ex}%[1D3N2-2]%[ma:75c3e073]
	 Giới hạn $\lim \limits_{x \to 0} \dfrac{6 - 2 x}{9 x - 4}$ bằng
	\choice 
		{ \True $-  \dfrac{3}{2}$}
		{  $\dfrac{3}{2}$}
		{ $\dfrac{2}{3}$} 
		{ $- \dfrac{2}{9}$} 
	\loigiai{ 
		Ta có  $\lim \limits_{x \to 0} \dfrac{6 - 2 x}{9 x - 4}=\dfrac{6-2\cdot 0}{9\cdot 0-4}=-\dfrac{3}{2}$.
	} 
\end{ex}

% nguồn: Ok_TracNghiemLop11_Lan1
\begin{ex}%[1D3N2-2]%[ma:5c9c2944]
	Giá trị của giới hạn $\lim\limits_{x\to -1}\dfrac{x^2-3}{x^3+2}$ là 
	 \choice
		 {$1$}
		 {\True $-2$}
		 {$2$}
		 {$-\dfrac{3}{2}$}
	 \loigiai{
		Ta có $\lim\limits_{x\to-1}\dfrac{x^2-3}{x^3+2}= \dfrac{(-1)^2-3}{(-1)^3+2}=-2$. 
	 }
\end{ex}

% nguồn: Ok_TracNghiemLop11_Lan1
\begin{ex}%[1D3N2-2]%[ma:96c55ff5]
	Giá trị của giới  hạn $\lim\limits_{x\to 0}{x^2}\sin\dfrac{1}{2}$ là
	\choice 
		 {$\sin\dfrac{1}{2}$}
		{$+\infty$} 
		 {$-\infty$}
		 {\True $0$}
	 \loigiai{
		Ta có $\lim\limits_{x\to 0}{x^2}\sin\dfrac{1}{2}=0 \cdot \sin\dfrac{1}{2}=0$. 
	 }
\end{ex}

% nguồn: Ok_TracNghiemLop11_Lan1
\begin{ex}%[1D3N2-2]%[ma:ff16a3cb]
	Giá trị của giới  hạn $\lim\limits_{x\to -1}\dfrac{|x-1|}{x^4+x-3}$ là
	 \choice
		 {$-\dfrac{3}{2}$}
		{$\dfrac{2}{3}$} 
		{$\dfrac{3}{2}$} 
		{\True  $-\dfrac{2}{3}$}
	 \loigiai{
		Ta có $\lim\limits_{x\to-1}\dfrac{|x-1|}{x^4+x-3}=  \dfrac{|-1-1|}{1-1-3}=-\dfrac{2}{3}$.
	 }
\end{ex}

% nguồn: Ok_TracNghiemLop11_Lan1
\begin{ex}%[1D3N2-2]%[ma:53a52298]
	Giá trị của $\lim\limits_{x\rightarrow -2} \left(x^3-x^2+1\right)$ bằng 
	 \choice
		 {\True $-11$}
		 {$12$}
		{$5$} 
		{$0$} 
	\loigiai{ 
		Ta có $\lim\limits_{x\rightarrow -2}  \left(x^3-x^2+1\right)=\lim\limits_{x\rightarrow -2} \left((-2)^3-(-2)^2+1\right)=-11$.
	} 
\end{ex}

% nguồn: Ok_TracNghiemLop11_Lan1
\begin{ex}%[1D3N2-2]%[ma:389d4665]
	Cho hai hàm số $f(x)$, $g(x)$ có giới hạn hữu hạn  tại $x=a$ đồng thời thỏa mãn các điều kiện\break $\lim\limits_{x\to a} [2f(x)-3g(x)]=3$ và $\lim\limits_{x\to a} [f(x)+6g(x)]=4$. Tính giới hạn $L=\lim\limits_{x\to a} [2(f(x)+g(x))]$.
	\choice 
		 {$L=\dfrac{7}{3}$}
		{$L=\dfrac{7}{6}$} 
		{\True  $L=\dfrac{14}{3}$}
		{$L=7$} 
	 \loigiai{
	 Ta có
	\[\heva{&\lim\limits_{x\to a} [2f(x)-3g(x)]=3 \\& \lim\limits_{x\to a} [f(x)+6g(x)]=4}\Leftrightarrow \heva{&2\lim\limits_{x\to a} f(x)-3\lim\limits_{x\to a} g(x)=3 \\& \lim\limits_{x\to a}f(x)+6\lim\limits_{x\to a}g(x) = 4 }\Leftrightarrow \heva{&\lim\limits_{x\to a}f(x)= 2\\& \lim\limits_{x\to a}g(x)=  \dfrac{1}{3}. }\]
	Vậy $L=\lim\limits_{x\to a} [2(f(x)+g(x))]=2\left[\lim\limits_{x\to a}f(x)+\lim\limits_{x\to a}g(x)\right] =  2\cdot\left(2+\dfrac{1}{3}\right) = \dfrac{14}{3}$.
	} 
\end{ex}

% nguồn: Ok_TracNghiemLop11_Lan1
\begin{ex}%[1D3N2-2]%[ma:948396de]
	Giá trị của $\lim \limits_{x  \rightarrow-4} \dfrac{x^2+3 x-4}{x^2+4 x}$ bằng
	\choice 
		{ $1$  }
		{$-1$  }
		{\True $\dfrac{5}{4}$} 
		{$-\dfrac{5}{4}$} 
	\loigiai{Ta có 
	$\lim  \limits_{x \rightarrow-4} \dfrac{x^2+3 x-4}{x^2+4 x}=\lim \limits_{x \rightarrow-4} \dfrac{(x-1)(x+4)}{x(x+4)}=\lim \limits_{x \rightarrow-4} \dfrac{(x-1)}{x}=\dfrac{5}{4} $. }
\end{ex}

% nguồn: Ok_TracNghiemLop11_Lan1
\begin{ex}%[1D3N2-2]%[ma:d21b30d0]
	Giới  hạn $\lim\limits_{x \to -2}{(2x+1)}$ có giá trị bằng
	 \choice
		 {$5$}
		 {\True $-3$}
		{$+ \infty  $}
		 {$- \infty $}
	 \loigiai{
		Ta có $\lim\limits_{x \to -2}{(2x+1)}$ = $2 \cdot (-2) + 1$ =  $-3$.
	} 
\end{ex}

% nguồn: Ok_TracNghiemLop11_Lan1
\begin{ex}%[1D3N2-2]%[ma:b0c21869]
	Giá trị của giới hạn $\lim\limits_{x\to 1}\dfrac{x-x^3}{(2x-1)(x^4-3)}$ là 
	 \choice
		 {$1$}
		{$-2$} 
		{\True $0$} 
		{$-\dfrac{3}{2}$} 
	\loigiai{ 
		Ta có $\lim\limits_{x\to 1}\dfrac{x-x^3}{(2x-1)(x^4-3)}= 
		\dfrac{1-1^3}{(2  \cdot 1-1)(1^4-3)}=0$.
	} 
\end{ex}

% nguồn: Ok_TracNghiemLop11_Lan1
\begin{ex}%[1D3N2-2]%[ma:89f39864]
	 Tính $\lim\limits_{x\to -2}(x^2+2x+2)$.
	 \choice
		{\True $2$} 
		{$10$} 
		{$-6$} 
		{$+\infty$} 
	\loigiai{ 
		Ta có $\lim\limits_{x\to -2} (x^2+2x+2)=(-2)^2+2\cdot  (-2)+2=2$.
	 }
\end{ex}

% nguồn: Ok_TracNghiemLop11_Lan1
\begin{ex}%[1D3N2-2]%[ma:63ab92cc]
	Cho giới hạn $\lim\limits_{x\to-2}\left(x^2-2 a  x+3+a^2\right)=3$ thì $a$ bằng bao nhiêu.
	 \choice {$a=2$}{$a=0$}{\True $a=-2$}{$a=-1$}
	 \loigiai{Ta có\\
		$\lim\limits_{x\to-2}\left(x^2-2 a  x+3+a^2\right)=3 \Leftrightarrow(-2)^2-2 a(-2)+3+a^2=3$ $\Leftrightarrow a^2+4 a+4=0 \Leftrightarrow a=-2$.
	 }
\end{ex}

% nguồn: Ok_Toan11_TracNghiem_Lan2
\begin{ex}%[1D3H2-2]%[ma:33f8ba25]
	Cho các giới hạn $\lim\limits_{x\to x_0}f(x)=2$; $\lim\limits_{x\to x_0}g(x)=3$, hỏi  $\lim\limits_{x\to x_0}[3f(x)-4g(x)]$ bằng
	\choice 
		 {$5$}
		{$2$} 
		{\True $-6$} 
		 {$3$}
	 \loigiai{
		Ta có $\lim\limits_{x\to x_0}[3f(x)-4 g(x)]=\lim\limits_{x\to x_0}3f(x)-\lim\limits_{x\to x_0}4g(x)=3 \lim\limits_{x\to x_0}f(x)-4 \lim\limits_{x\to  x_0}g(x)=3\cdot 2-4\cdot 3=-6$.
	 }
\end{ex}

% nguồn: Ok_Toan11_TracNghiem_Lan2
\begin{ex}%[1D3H2-2]%[ma:6f2b5bd5]
	Tìm $m$ để $A=\dfrac{1}{2}$ với $A=\lim\limits_{x\to  1}\dfrac{4x+m}{mx-2}$
	 \choice
		{$3$} 
		 {$2$}
		{\True $-10$} 
		 {$\dfrac{10}{3}$}
	\loigiai{ 
	 Ta có $A=\lim\limits_{x\to 1}\dfrac{4x+m}{mx-2}=\dfrac{4+m}{m-2}=\dfrac{1}{2}\Rightarrow m=-10$.}
\end{ex}

% nguồn: Ok_Toan11_TracNghiem_Lan2
\begin{ex}%[1D3H2-2]%[ma:0dd9904f]
	$\lim\limits_{x\to \sqrt{3}}\left|x^2-4\right|$  bằng
	 \choice
		{$-5$} 
		{\True $1$} 
		 {$5$}
		{$-1$} 
	 \loigiai{
		Ta có  $\lim\limits_{x\to \sqrt{3}}\left|x^2-4\right|=|3-4|=1$.
	} 
\end{ex}

% nguồn: Ok_Toan11_TracNghiem_Lan2
\begin{ex}%[1D3H2-2]%[ma:3cc5603d]
	 ]
	Giá trị của  $\lim\limits_{x\to 1}\left(2x^2-3x+1\right)$ bằng
	 \choice
		 {$2$}
		{$1$} 
		 {$+\infty$}
		{\True  $0$}
	 \loigiai{
		 Ta có $\lim\limits_{x\to 1}\left(2x^2-3x+1\right)=2
		 cdot 1^2+3
		 cdot 1+1=0$.
	 }
\end{ex}

% nguồn: Ok_Toan11_TracNghiem_Lan2
\begin{ex}%[1D3H2-2]%[ma:300cf9f8]
	Tính $\lim\limits_{x\to  1}\dfrac{x^3-2x^2+2020}{2x-1}$.
	\choice 
		 {$0$}
		 {$-\infty$}
		 {$+\infty$}
		{\True $2019$} 
	\loigiai{ 
		Ta  có $\lim\limits_{x\to 1}\dfrac{x^3-2x^2+2020}{2x-1}=\dfrac{1^3-2\cdot 1^2+2020}{2\cdot 1-1}=2019$.
	 }
\end{ex}

% nguồn: Ok_Toan11_TracNghiem_Lan2
\begin{ex}%[1D3H2-2]%[ma:e74bd29a]
	Giới hạn  $\lim\limits_{x\to 1}\dfrac{x^2-2 x+3}{x+1}$ bằng
	 \choice
		{\True $1$} 
		 {$0$}
		{$3$} 
		{$2$} 
	\loigiai{ 
		Ta có $\lim\limits_{x\to  1}\dfrac{x^2-2 x+3}{x+1}=\dfrac{1^2-2\cdot 1+3}{1+1}=1$.
	 }
\end{ex}

% nguồn: Ok_Toan11_TracNghiem_Lan2
\begin{ex}%[1D3H2-2]%[ma:a2f52221]
	Giới hạn nào sau đây có kết quả bằng $+\infty$? 
	 \choice
		{$\lim\limits_{x\to 1}\dfrac{x-3}{(x-1)^2}$} 
		 {$\lim\limits_{x\to 1}\dfrac{x-2}{(x-1)^2}$}
		 {$\lim\limits_{x\to 1}\dfrac{-x-1}{(x-1)^2}$}
		 {\True $\lim\limits_{x\to 1}\dfrac{x+1}{(x-1)^2}$}
	 \loigiai{
		Ta có $(x-1)^2\geq 0$, $\forall x\neq 1$, suy ra để giới hạn bằng $+\infty$ thì giới hạn của tử phải  dương khi $x\to 1$.\\
		Vậy  $\lim\limits_{x\to 1}\dfrac{x+1}{(x-1)^2}=+\infty$.
	} 
\end{ex}

% nguồn: Ok_Toan11_TracNghiem_Lan2
\begin{ex}%[1D3H2-2]%[ma:cb1f592e]
	Giá trị của $\lim\limits_{x \to 1}\dfrac{x^{2018}+x-2}{x^{2017}+x-2}$ bằng $\dfrac{a}{b}$, với $\dfrac{a}{b}$ là phân số tối giản.  Tính giá trị của $a^2-b^2$.
	 \choice
		{\True  $4\,037$}
		 {$4\,035$}
		{$-4\,035$} 
		 {$4\,033$}
	 \loigiai{
		Ta  có
		\begin{eqnarray*}
			&&\lim\limits_{x \to 1}\dfrac{x^{2018}+x-2}{x^{2017}+x-2}\\ 
			&=&\lim\limits_{x  \to 1}\dfrac{x^{2018}-1+x-1}{x^{2017}-1+x-1}\\
			&=&\lim\limits_{x \to 1}\dfrac{\left( x-1 \right)\left( x^{2017}+x^{2016}+\ldots+x+1 \right)+x-1}{\left( x-1 \right)\left( x^{2016}+x^{2015}+\ldots+x+1 \right)+x-1}\\ 
			&=&\lim\limits_{x  \to 1} \dfrac{x^{2017}+x^{2016}\ldots+x+2}{x^{2016}+x^{2015}+\ldots+x+2}\\
			 &=&\dfrac{1+1+\ldots+1+2}{1+1+\ldots+1+2}=\dfrac{2\,019}{2\,018}.
		\end{eqnarray*}
	Vậy  $a^2-b^2=4\,037$.}
\end{ex}

% nguồn: Ok_Toan11_TracNghiem_Lan2
\begin{ex}%[1D3H2-2]%[ma:dfc8aea3]
	Tính  giới hạn $\lim\limits_{x \to 1}\dfrac{x^2-3x+2}{x-1}$.
	\choice 
		{$1$} 
		{\True  $-1$}
		 {$2$}
		 {$-2$}
	\loigiai{ 
	Ta có $\lim\limits_{x \to 1}\dfrac{x^2-3x+2}{x-1}=\lim\limits_{x \to 1}\dfrac{(x-1)(x-2)}{x-1}=\lim\limits_{x \to  1}(x-2)=-1$.}
\end{ex}

% nguồn: Ok_Toan11_TracNghiem_Lan2
\begin{ex}%[1D3H2-2]%[ma:e07b98c7]
	Tìm giới hạn  $A=\lim\limits_{x\to-2}\dfrac{x+1}{x^2+x+4}$.
	 \choice
		{\True  $-\dfrac{1}{6}$}
		{$-\infty$} 
		{$+\infty$} 
		 {$1$}
	\loigiai{ 
		Ta có $A=\lim\limits_{x\to-2}\dfrac{x+1}{x^2+x+4}=\dfrac{(-2)+1}{(-2)^2+(-2)+4}=-\dfrac{1}{6}$. 
	} 
\end{ex}

% nguồn: Ok_Toan11_TracNghiem_Lan2
\begin{ex}%[1D3H2-2]%[ma:66338957]
	Cho giới hạn $\lim\limits_{x \to 2}\dfrac{x^2-3x+2}{x^2-4}=\dfrac{a}{b}$ trong  đó $\dfrac{a}{b}$ là phân số tối giản. Tính $S=a^2+b^2$.
	 \choice
		 {$S=20$}
		{\True  $S=17$}
		{$S=10$} 
		 {$S=25$}
	 \loigiai{
		$\lim\limits_{x \to  2}\dfrac{x^2-3x+2}{x^2-4}=\lim\limits_{x \to 2}\dfrac{(x-1)(x-2)}{(x+2)(x-2)}=\lim\limits_{x \to 2}\dfrac{x-1}{x+2}=\dfrac{1}{4}$.
	Do đó $a=1$, $b=4$ suy ra $S=1^2+4^2=17$.} 
\end{ex}

% nguồn: Ok_Toan11_TracNghiem_Lan2
\begin{ex}%[1D3H2-2]%[ma:3f890ae2]
	Tính giới hạn $I=\lim\limits_{x \to 2}\dfrac{x^2-5x+6}{x-2}$. 
	 \choice
		 {\True $I=-1$}
		{$I=0$} 
		 {$I=1$}
		{$I=5$} 
	 \loigiai{
	Ta có $I=\lim\limits_{x \to 2}\dfrac{x^2-5x+6}{x-2}=\lim\limits_{x \to 2}\dfrac{\left( x-2 \right)\left( x-3 \right)}{x-2}=\lim\limits_{x \to 2}\left( x-3 \right)=-1$.} 
\end{ex}

% nguồn: Ok_Toan11_TracNghiem_Lan2
\begin{ex}%[1D3H2-2]%[ma:45782b46]
	 Cho $\lim\limits_{x\to 3}f(x)=-2$. Tính $\lim\limits_{x\to 3}[f(x)+4x-1]$.
	 \choice
		 {$5$}
		 {$6$}
		 {$11$}
		 {\True $9$}
	\loigiai{ 
		Ta có $\lim\limits_{x\to 3}[f(x)+4x-1]=-2+4\cdot 3-1=9$. 
	} 
\end{ex}

% nguồn: Ok_Toan11_TracNghiem_Lan2
\begin{ex}%[1D3H2-2]%[ma:40cf4483]
	Tìm $m$ để $A=7$ với  $A=\lim\limits_{x\to 1}\dfrac{4x+m^2-1}{-x+2}$
	 \choice
		 {\True $\pm 2$}
		{$2$} 
		{$-2$} 
		{$\dfrac{10}{3}$} 
	 \loigiai{
		Ta có $A=\lim\limits_{x\to 1}\dfrac{4x+m^2-1}{-x+2}=\dfrac{4+m^2-1}{-1+2}=3+m^2=7\Rightarrow  m=\pm 2$.
	 }
\end{ex}

% nguồn: Ok_Toan11_TracNghiem_Lan2
\begin{ex}%[1D3H2-2]%[ma:98744a9a]
	Giới  hạn $\lim\limits_{x \to -2}\dfrac{x+1}{{{\left( x+2 \right)}^{2}}}$ bằng
	 \choice
		 {\True $-\infty $}
		 {$\dfrac{3}{16}$}
		 {$0$}
		{$+\infty $} 
	\loigiai{ 
		Ta có $\lim\limits_{x \to  -2}\dfrac{x+1}{{{\left( x+2 \right)}^{2}}}=\lim\limits_{x \to -2}\dfrac{1}{{{\left( x+2 \right)}^{2}}}\cdot \left( x+1 \right)=-\infty $.\\
	Do $\lim\limits_{x \to -2}\dfrac{1}{{{\left( x+2 \right)}^{2}}}=+\infty $ và $\lim\limits_{x  \to -2}\left( x+1 \right)=-1<0$.}
\end{ex}

% nguồn: Ok_Toan11_TracNghiem_Lan2
\begin{ex}%[1D3H2-2]%[ma:0a75bd93]
	Tính $\lim\limits_{x \to 3}\dfrac{x^2-9}{x-3}$ bằng 
	 \choice
		{$3$} 
		 {\True $6$}
		{$+\infty  $}
		{$-3$} 
	\loigiai{ 
	 Ta có $\lim\limits_{x \to 3}\dfrac{x^2-9}{x-3}=\lim\limits_{x \to 3}\left( x+3 \right)=6$.}
\end{ex}

% nguồn: Ok_Toan11_TracNghiem_Lan2
\begin{ex}%[1D3H2-2]%[ma:92619ae0]
	Tính  giới hạn $A=\lim\limits_{x \to 1}\dfrac{x^3-1}{x-1}$.
	\choice 
		{$A=-\infty$} 
		 {$A=0$}
		 {\True $A=3$}
		 {$A=+\infty$}
	\loigiai{ 
	$A=\lim\limits_{x \to 1}\dfrac{x^3-1}{x-1}$$=\lim\limits_{x \to 1}\dfrac{\left( x-1 \right)\left( x^2+x+1 \right)}{x-1}$$=\lim\limits_{x \to 1}\left(  x^2+x+1 \right)=3$.}
\end{ex}

% nguồn: Ok_Toan11_TracNghiem_Lan2
\begin{ex}%[1D3H2-2]%[ma:7b982ab4]
	Với giá trị nào của tham số $m$ thì  $\lim\limits_{x\to -1}\left( mx^2+3x-2m \right)=0$?
	 \choice
		 {\True $m=-3$}
		 {$m=-1$}
		{$m=0$} 
		{$m=3$} 
	 \loigiai{
	Ta có $\lim\limits_{x\to -1}\left(  mx^2+3x-2m \right)=0\Leftrightarrow m\cdot {{\left( -1 \right)}^2}+3\cdot \left( -1 \right)-2m=0\Leftrightarrow m=-3$.}
\end{ex}

% nguồn: Ok_Toan11_TracNghiem_Lan2
\begin{ex}%[1D3H2-2]%[ma:51b58535]
	Biết $\lim\limits_{x\to 1}\dfrac{x+\sqrt{x^2+1}}{x+1}=a+b\sqrt{2}$. Tính $a+b$  được kết quả đúng bằng
	\choice 
		{$2$} 
		 {\True $1$}
		{$5$} 
		{$0$} 
	\loigiai{ 
	 Ta có $\lim\limits_{x\to 1}\dfrac{x+\sqrt{x^2+1}}{x+1}=\dfrac{1+\sqrt{2}}{2}=\dfrac{1}{2}+\dfrac{1}{2}\sqrt{2}\Rightarrow a=\dfrac{1}{2},b=\dfrac{1}{2}\Rightarrow a+b=1$.}
\end{ex}

% nguồn: Ok_Toan11_TracNghiem_Lan2
\begin{ex}%[1D3H2-2]%[ma:f1be54ce]
	Giá  trị của giới hạn $\lim\limits_{x\to 2}(3x^2+7x+11)$ là
	\choice 
		{\True  $37$}
		{$38$} 
		{$39$} 
		{$40$} 
	\loigiai{ 
		Ta có $\lim\limits_{x\to 2}(3x^2+7x+11)=3\cdot 2^2+7\cdot 2+11=37$. 
	 }
\end{ex}

% nguồn: Ok_Toan11_TracNghiem_Lan2
\begin{ex}%[1D3H2-2]%[ma:585d7cb5]
	Tìm $m$ để  $A=3$ với $A=\lim\limits_{x\to 2}\dfrac{3x+m}{x+2}$
	 \choice
		{\True  $6$}
		{$14$} 
		{$3$} 
		 {$\dfrac{10}{3}$}
	 \loigiai{
	 Ta có $A=\lim\limits_{x\to 2}\dfrac{3x+m}{x+2}$ $=\dfrac{3\cdot 2+m}{2+2}=\dfrac{6+m}{4}=3\Rightarrow m=6$.}
\end{ex}

% nguồn: Ok_Toan11_TracNghiem_Lan2
\begin{ex}%[1D3H2-2]%[ma:3e2db213]
	 Giá trị của giới hạn $\lim\limits_{x\to \sqrt{3}}\left| x^2-4\right|$ là
	\choice 
		{\True $1$} 
		 {$-1$}
		{$2$} 
		 {$3$}
	\loigiai{ 
		Ta có $\lim\limits_{x\to  \sqrt{3}}\left| x^2-4\right|=\left|\sqrt3^2-4\right|=1$.
	 }
\end{ex}

% nguồn: Ok_Toan11_TracNghiem_Lan2
\begin{ex}%[1D3H2-2]%[ma:da351ee1]
	$\lim\limits_{x\to-2}\dfrac{2|x+1|-5  \sqrt{x^2-3}}{2 x+3}$ bằng
	\choice 
		{$\dfrac{1}{3}$} 
		 {$\dfrac{1}{7}$}
		 {$7$}
		 {\True $3$}
	 \loigiai{
		Ta  có $\lim\limits_{x\to-2}\dfrac{2|x+1|-5\sqrt{x^2-3}}{2x+3}=\dfrac{2-5}{-1}=3$.
	 }
\end{ex}

% nguồn: Ok_Toan11_TracNghiem_Lan2
\begin{ex}%[1D3H2-2]%[ma:db18f945]
	Giá trị của  $\lim\limits_{x\to 1}\left(3 x^2-2 x+1\right)$ bằng
	\choice 
		 {$+\infty$}
		{\True $2$} 
		 {$1$}
		 {$3$}
	\loigiai{ 
		 Ta có $\lim\limits_{x\to 1}\left(3 x^2-2 x+1\right)=3\cdot 1^2-2\cdot 1+1=2$.
	} 
\end{ex}

% nguồn: Ok_Toan11_TracNghiem_Lan2
\begin{ex}%[1D3H2-2]%[ma:a7c90215]
	$\lim\limits_{x\to 5^+} \dfrac{\left| 10-2x \right|}{x^2-6x+5}$ là 
	 \choice
		 {$+\infty $}
		{$0$} 
		{$-\dfrac{1}{2}$} 
		{\True  $\dfrac{1}{2}$}
	\loigiai{ 
	$\lim\limits_{x \to {{5}^{+}}}\dfrac{\left| 10-2x  \right|}{x^2-6x+5}=\lim\limits_{x \to {{5}^{+}}}\dfrac{2x-10}{x^2-6x+5}=\lim\limits_{x \to {{5}^{+}}}\dfrac{2}{x-1}=\dfrac{1}{2}$.}
\end{ex}

% nguồn: Ok_Toan11_TracNghiem_Lan2
\begin{ex}%[1D3H2-2]%[ma:1c4f3852]
	Giới  hạn $\lim\limits_{x\to-1}\left(x^2-x+7\right)$ bằng
	 \choice
		{$5$} 
		{\True $9$} 
		{$0$} 
		 {$7$}
	\loigiai{ 
		Ta có $\lim\limits_{x\to-1}\left(x^2-x+7\right)=(-1)^2-(-1)+7=9$. 
	 }
\end{ex}

% nguồn: Ok_Toan11_TracNghiem_Lan2
\begin{ex}%[1D3H2-2]%[ma:d3cd8d7f]
	Biết  $\lim\limits_{x\to 1}\dfrac{x^2-ax+1}{x+1}=3$. Khi đó giá trị của $a$ là
	 \choice
		 {$4$}
		{$0$} 
		 {\True $-4$}
		{$3$} 
	\loigiai{ 
	Có $\lim\limits_{x\to 1}\dfrac{x^2-ax+1}{x+1}=3\Rightarrow \dfrac{1^2-a\cdot  1+1}{1+1}=3\Rightarrow a=-4$.}
\end{ex}

% nguồn: Ok_Toan11_TracNghiem_Lan2
\begin{ex}%[1D3H2-2]%[ma:776a7e99]
	 $\lim\limits_{x\to 1}\dfrac{x+1}{x+2}$ bằng
	\choice 
		{$+\infty$} 
		 {$\dfrac{1}{2}$}
		 {\True $\dfrac{2}{3}$}
		 {$-\infty$}
	 \loigiai{
		Ta  có $\lim\limits_{x\to 1}\dfrac{x+1}{x+2}=\dfrac{1+1}{1+2}=\dfrac{2}{3}$.
	 }
\end{ex}

% nguồn: Ok_Toan11_TracNghiem_Lan2
\begin{ex}%[1D3H2-2]%[ma:f2d85313]
	Tính  $\lim\limits_{x \to 5}\dfrac{x^2-12x+35}{25-5x}$.
	 \choice
		 {$-\dfrac{2}{5}$}
		 {$+\infty $}
		{\True $\dfrac{2}{5}$} 
		{$-\infty  $}
	\loigiai{ 
	Ta có $\lim\limits_{x  \to 5}\dfrac{x^2-12x+35}{25-5x}=\lim\limits_{x \to 5}\dfrac{\left( x-7 \right)\left( x-5 \right)}{-5\left( x-5 \right)}=\lim\limits_{x \to 5}\dfrac{x-7}{-5}=\dfrac{2}{5}$.}
\end{ex}

% nguồn: Ok_Toan11_TracNghiem_Lan2
\begin{ex}%[1D3H2-2]%[ma:bc65222f]
	Kết quả của giới hạn $\lim\limits_{x  \to 2}\dfrac{x^2-4}{x-2}$ bằng
	 \choice
		{$0$} 
		 {\True $4$}
		 {$-4$}
		{$2$} 
	\loigiai{ 
	Ta có $\lim\limits_{x \to  2}\dfrac{x^2-4}{x-2}=\lim\limits_{x \to 2}\dfrac{\left( x-2 \right)\left( x+2 \right)}{x-2}=\lim\limits_{x \to 2}\left( x+2 \right)=4$.}
\end{ex}

% nguồn: Ok_TracNghiemLop11_Lan1
\begin{ex}%[1D3H2-2]%[ma:23b5a15f]
	Giới hạn  $\lim\limits_{x\to 1}\, \dfrac{x-3}{x+2}$ bằng
	\choice 
		{$0$} 
		{\True  $\dfrac{-2}{3}$}
		{$+\infty  $}
		 {$1$}
	 \loigiai{
		Ta có $\lim\limits_{x\to  1}\, \dfrac{x-3}{x+2}=\dfrac{1-3}{1+2}=\dfrac{-2}{2}$.
	} 
\end{ex}

% nguồn: Ok_TracNghiemLop11_Lan1
\begin{ex}%[1D3H2-2]%[ma:106e6455]
	 Giá trị của giới hạn $\lim\limits_{x\to 2}\dfrac{\sqrt[3]{3x^2-4}-\sqrt{3x-2}}{x+1}$ là
	\choice 
		 {$-\dfrac{3}{2}$}
		 {$-\dfrac{2}{3}$}
		{\True $0$} 
		 {$+\infty$}
	\loigiai{ 
		Ta có $\lim\limits_{x\to 2}\dfrac{\sqrt[3]{3x^2-4}-\sqrt{3x-2}}{x+1}= \dfrac{\sqrt[3]{12-4}-\sqrt{6-2}}{3}=\dfrac{0}{3}=0$. 
	} 
\end{ex}

% nguồn: Ok_TracNghiemLop11_Lan1
\begin{ex}%[1D3K2-2]%[ma:03d6427a]
	Cho dãy số các định bởi $\left\{\begin{aligned}
	&u_1=2018\\ 
	&u_{n+1}=\sqrt{u^2_n+n^2+2018};\,  n\geq 1
\end{aligned}\right.$. Số hạng thứ $21$ trong dãy số có giá trị gần nhất là
 \choice
	{$201$} 
	{$207$} 
	 {\True$213$}
	{$219$} 
 \loigiai{
Ta  có
 {\allowdisplaybreaks
\begin{eqnarray*}
	&u_{n+1}& = \sqrt{u^2_n+n^2+2018} \Rightarrow u^2_{n+1}=u^2_n+n^2+2018;\,  n\geq 1\\
	&u^2_1&=u^2_1\\ 
	 &u^2_2&=u^2_2+1^2+2018\\
	 &u^2_3&=u^2_2+2^2+2018\\
	&u^2_4&=u^2_3+3^2+2018\\ 
	 &\ldots& \quad \ldots\\
	&u^2_n&=u^2_{n-1}+(n-1)^2+2018\\ 
\end{eqnarray*}
} 
Cộng  $n$ đẳng thức trên theo vế ta được
\[u^2_n=1^2+2^2+3^2+\cdots+(n-1)^2+2018n\] 
 Mà
 \[1^2+2^2+3^2+\cdots+n^2=\dfrac{n(n+1)(2n+1)}{6}\]
 Do đó
\[1^2+2^2+3^2+\cdots+(n-1)^2=\dfrac{(n-1)n(2n-1)}{6}\] 
 Nên
 \[u_n=\dfrac{1}{6}\sqrt{6n(2n^2-3n+12109)}\Rightarrow u_{21}=8 \sqrt{707}\approx 213\]
} 
\end{ex}

% nguồn: Ok_TracNghiemLop11_Lan1
\begin{ex}%[1D3Y2-2]%[ma:bdc1f4f5]
	Cho dãy  số $({{u}_{n}}),$ biết ${{u}_{n}}={{(-1)}^{n}}\cdot \dfrac{{{2}^{n}}}{n}$. Tìm số hạng ${{u}_{3}}$.
	 \choice
		{  \True ${{u}_{3}}=-\dfrac{8}{3}$}
		{${{u}_{3}}=2$ } 
		 {${{u}_{3}}=-2$ }
		{${{u}_{3}}=\dfrac{8}{3}$ } 
	 \loigiai{
		Thay trực tiếp hoặc dùng chức  năng CALC:\\
	$${{u}_{3}}={{(-1)}^3}\cdot  \dfrac{{{2}^3}}{3}=-\dfrac{8}{3}.$$}
\end{ex}

% nguồn: Ok_TracNghiemLop11_Lan1
\begin{ex}%[1D3B2-2]%[ma:fcc81d20]
	Cho dãy số $(u_n)$ có số hạng tổng quát $u_n=\dfrac{2n-1}{n+2}$. Tìm số hạng thứ  $100$ và $200$
	\choice 
		{$u_{100}=\dfrac{2\cdot100-1}{100+2}=\dfrac{103}{34}$; $u_{200}=\dfrac{2\cdot200+1}{200+2}=\dfrac{309}{202}$} 
		 {$u_{100}=\dfrac{2\cdot100-1}{100+2}=\dfrac{130}{34}$; $u_{200}=\dfrac{2\cdot200+1}{200+2}=\dfrac{390}{202}$}
		{$u_{100}=\dfrac{2\cdot100-1}{100+2}=\dfrac{33}{34}$; $u_{200}=\dfrac{2\cdot200+1}{200+2}=\dfrac{392}{202}$} 
		{\True$u_{100}=\dfrac{2\cdot100-1}{100+2}=\dfrac{133}{34}$; $u_{200}=\dfrac{2\cdot200+1}{200+2}=\dfrac{399}{202}$} 
	 \loigiai{
		Số  hạng thứ 100: $u_{100}=\dfrac{2\cdot100-1}{100+2}=\dfrac{133}{34}$.\\
	Số hạng thứ  200: $u_{200}=\dfrac{2\cdot200-1}{200+2}=\dfrac{399}{202}$.}
\end{ex}

% nguồn: Ok_TracNghiemLop11_Lan1
\begin{ex}%[1D3B2-2]%[ma:de8142a4]
	Cho dãy số $\left(a_n\right)$ có $a_n=-n^2+4n+11,\forall n\in \mathbb{N}^*$. Tìm số hạng  lớn nhất của dãy số $\left(a_n\right)$.
	 \choice
		 {$14$}
		 {\True $15$}
		 {$13$}
		{$12$} 
	\loigiai{ 
		 Ta có $a_n=-(n-2)^2+15\le 15,\forall n\ge 1.$ Dấu bằng xảy ra khi $n-2=0\Leftrightarrow n=2.$\\
	Vậy số hạng lớn nhất của dãy số là số hạng bằng 15} 
\end{ex}

% nguồn: Ok_TracNghiemLop11_Lan1
\begin{ex}%[1D3B2-2]%[ma:2dc3f422]
	Cho dãy số $\left(a_n\right)$ có $a_n=\dfrac{n}{n^2+100},\forall n\in \mathbb{N}^*$. Tìm số hạng lớn nhất của dãy số $\left(a_n\right)$. 
	\choice 
		 {\True $\dfrac{1}{20}$}
		{$\dfrac{1}{30}$} 
		 {$\dfrac{1}{25}$}
		{$\dfrac{1}{21}$} 
	\loigiai{ 
		Ta có $a_n=\dfrac{n}{n^2+100}\le \dfrac{n}{2\sqrt{n^2.100}}=\dfrac{1}{20}.$ Dấu  bằng xảy ra khi $n^2=100\Leftrightarrow n=10.$\\
	Vậy số hạng lớn nhất của dãy  là số hạng bằng $\dfrac{1}{20}$}
\end{ex}

% nguồn: Ok_TracNghiemLop11_Lan1
\begin{ex}%[1D3B2-2]%[ma:792553b2]
	Cho dãy số $(u_n)$ có số hạng tổng quát $u_n=\dfrac{2n+1}{n+2}$. Số $\dfrac{167}{84}$  là số hạng thứ mấy?
	 \choice
		{300} 
		 {212}
		 {\True 250}
		{249} 
	 \loigiai{
		Giả sử $u_n=\dfrac{167}{84}\Rightarrow  \dfrac{2n+1}{n+2}=\dfrac{167}{84}\Leftrightarrow 84(2n+1)=167(n+2)$$\Leftrightarrow n=250$.\\
	Vậy $\dfrac{167}{84}$ là số hạng thứ $250$  của dãy số $(u_n)$.}
\end{ex}

% nguồn: Ok_TracNghiemLop11_Lan1
\begin{ex}%[1D3K2-2]%[ma:b8e350e5]
	Cho dãy số $\left(x_n\right)$ xác định bởi $x_1=\dfrac{2}{3}$ và ${x}_{n+1}=\dfrac{x_n}{2\left(2n+1\right)x_n+1},\forall n\in \mathbb{N}^*$. Mệnh đề nào  dưới đây là đúng ?
	 \choice
		 {\True$x_{100}=\dfrac{2}{39999}$}
		 {$x_{100}=\dfrac{39999}{2}$}
		 {$x_{100}=\dfrac{2}{40001}$}
		 {$x_{100}=\dfrac{2}{40803}$}
	\loigiai{ 
		Ta có $x_n>0,\forall  n\ge 1$ và $\dfrac{1}{{x}_{n+1}}=2(2n+1)+\dfrac{1}{x_n},\forall n\ge 1.$\\
		Suy  ra $\dfrac{1}{x_n}=\dfrac{1}{x_1}+4(1+2+\cdots+n-1)+2(n-1)=\dfrac{3}{2}+2n(n-1)+2(n-1)=\dfrac{4n^2-1}{2}.$\\
	 Suy ra $x_n=\dfrac{2}{4n^2-1}.$ Do đó $x_{100}=\dfrac{2}{39999}.$}
\end{ex}

% nguồn: Ok_TracNghiemLop11_Lan1
\begin{ex}%[1D3B2-2]%[ma:1a355848]
	Cho dãy số $(u_n)$ có số hạng tổng quát:  $u_n = 2n+\sqrt{n^2+4}$. Dãy số đã cho có bao nhiêu số hạng là số nguyên.
	 \choice
		 {$1$}
		 {$2$}
		{$3$} 
		{\True$0$} 
	 \loigiai{
		Ta  có $u_n$ nguyên $\Leftrightarrow \sqrt{n^2+4}=k\in\mathbb{N} \Leftrightarrow k^2-n^2=4 \Leftrightarrow (k-n)(k+n)=4$. Phương trình này vô nghiệm.\\
	Vậy  không có số hạng nào của dãy nhận giá trị nguyên.}
\end{ex}

% nguồn: Ok_TracNghiemLop11_Lan1
\begin{ex}%[1D3B2-2]%[ma:be484a39]
	Cho dãy số có 4 số hạng đầu là $-1,3,19,53$. Hãy tìm một quy luật của dãy số trên và viết số hạng thứ 10 của dãy với quy luật  vừa tìm.
	\choice 
		{$u_{10}=97$} 
		 {$u_{10}=71$}
		 {$u_{10}=1414$}
		 {\True $u_{10}=971$}
	 \loigiai{
		Xét dãy $(u_n)$ có dạng: $u_n=an^3+bn^2+cn+d$.\\ 
		Ta  có hệ $\heva{
		&  a+b+c+d=-1 \\
		&  8a+4b+2c+d=3 \\
		& 27a+9b+3c+d=19 \\ 
		 & 64a+16b+4c+d=53. \\}$\\
		Giải hệ trên ta tìm được: $a=1,b=0,c=-3,d=1$$\Rightarrow  u_n=n^3-3n+1$ là một quy luật.\\
	Số hạng thứ 10: $u_{10}=971$.} 
\end{ex}

% nguồn: Ok_TracNghiemLop11_Lan1
\begin{ex}%[1D3B2-2]%[ma:e7ffcfc4]
	Cho  dãysố $\heva{
	&u_1=4 \\ 
	& {u}_{n+1}=u_n+n \\}$. Tìm số hạng  thứ $5$ của dãy số.
	\choice 
		{$16$} 
		{$12$} 
		 {$15$}
		{\True $14$} 
	\loigiai{Ta có $u_2=u_1+1=5$; $u_3=u_2+2=7$; $u_4=u_3+3=10$. Do đó số  hạng thứ $5$ của dãy số là $u_5=u_4+4=14$}
\end{ex}

% nguồn: Ok_TracNghiemLop11_Lan1
\begin{ex}%[1D3B2-2]%[ma:f0549335]
	Cho dãysố $\left(u_n\right)$ xác định bởi $u_1=\dfrac{1}{2}$ và  $u_n={u}_{n-1}+2n$ với mọi $n\ge 2$. Khi đó $u_{50}$ bằng
	 \choice
		{$1274,5$} 
		{$2548,5$} 
		{$5096,5$} 
		 {\True $2550,5$}
	\loigiai{ 
		 Ta có $u_n=\dfrac{1}{2}+2(1+2+\cdots+n)=\dfrac{1}{2}+n(n+1)$.\\
	 Suy ra $u_{50}=\dfrac{1}{2}+50.51=2550,5.$}
\end{ex}

% nguồn: Ok_TracNghiemLop11_Lan1
\begin{ex}%[1D3B2-2]%[ma:230aaa7b]
	Cho dãy số $(u_n)$ được xác định bởi $u_n=\dfrac{n^2+3n+7}{n+1}$. 
	Dãy số có bao  nhiêu số hạng nhận giá trị nguyên.
	 \choice
		{$2$} 
		 {$4$}
		{\True $1$} 
		 {Không có}
	 \loigiai{
		Ta có $u_n=n+2+\dfrac{5}{n+1}$, do đó $u_n$ nguyên khi và chỉ khi $\dfrac{5}{n+1}$ nguyên hay $n+1$ là ước của $5$. \\ 
		Điều đó xảy ra  khi $n+1=5\Leftrightarrow n=4$.\\
	Vậy dãy số có  duy nhất một số hạng nguyên là $u_4=7$}
\end{ex}

% nguồn: Ok_TracNghiemLop11_Lan1
\begin{ex}%[1D3G2-2]%[ma:5aee9a29]
	Cho dãy số $(u_{n})$ xác định bởi $\left\{\begin{aligned}
	&u_1=1\cdot2\cdot3\\ 
	 &u_2=2\cdot3\cdot4\\
	 &u_n=n(n+1)(n+2)
\end{aligned}\right.$. Đặt $S_n=a_1+a_2+\cdots+a_n$. Giá trị của $S_{30}$ là
\choice 
	 {$28184$}
	{\True$245520$} 
	 {$215760$}
	{$278256$} 
 \loigiai{
 Ta có
\begin{eqnarray*}
	&S_1&=u_1=1\cdot2\cdot3\\ 
	 &S_2&=u_1+u_2=2\cdot3\cdot4\\
	&S_3&=u_1+u_2+u_3=2\cdot3\cdot5+3\cdot4\cdot5=3\cdot5\cdot6 
\end{eqnarray*}
$\Rightarrow S_1=\dfrac{1}{4}\cdot1\cdot2\cdot3\cdot4$, $S_2=\dfrac{1}{4}\cdot2\cdot3\cdot4\cdot5$, $S_3=\dfrac{1}{4}\cdot3\cdot4\cdot5\cdot6$\\ 
Nhận thấy quy luật nên giả sử $S_k=\dfrac{1}{4}\cdot k\cdot(k+1)(k+2)(k+3)$, $k\geq  3$ (giả thiết quy nạp).\\
Ta sẽ chứng minh  $S_{k+1}=\dfrac{1}{4}\cdot(k+1)(k+2)(k+3)(k+4)$.\\
 Thật vậy, theo đề bài $\Rightarrow S_{k+1}=S_k +u_{k+1}=S_k +(k+1)(k+2)(k+3)$.\\
Theo giả  thiết quy nạp $\Rightarrow S_{k+1} =\dfrac{1}{4}\cdot k(k+1)(k+2)(k+3)+(k+1)(k+2)(k+3)$.\\
$\Rightarrow  S_{k+1}=\dfrac{1}{4}(k+1)(k+2)(k+3)(k+4)$.\\
Theo nguyên tắc quy nạp suy ra $S_n=\dfrac{1}{4}\cdot n(n+1)(n+2)(n+3)\Rightarrow  S_{30}=245520$.
 }
\end{ex}

% nguồn: Ok_TracNghiemLop11_Lan1
\begin{ex}%[1D3K2-2]%[ma:6ffb1beb]
	Cho dãy số $(a_n)$ xác định bởi $a_1=-3$ và  $a_{n + 1}=a_n + n^2 - 3n + 4,\forall n\in\mathbb{N}^{*}$. Số $1391$ là số hạng thứ mấy của dãy số đã cho?
	\choice 
		 {\True$18$}
		{$17$} 
		{$20$} 
		 {$19$}
	 \loigiai{
	Từ hệ thức truy hồi của  dãy số $(a_n)$ ta có
	 $$
	a_n=a_1 + \left[1^2  + 2^2 + \ldots + (n - 1)^2\right] - 3[1 + 2 + \ldots + (n - 1)] + 4(n - 1)\Leftrightarrow a_n=\dfrac{n^3 - 6n^2 + 17n - 21}{3}.$$
	Suy ra số hạng tổng quát của dãy số  $(a_n)$ là $a_n=\dfrac{n^3 - 6n^2 + 17n - 21}{3}$.
	 Giải phương trình $a_n=1391$ ta được $n=18$.
	 }
\end{ex}

% nguồn: Ok_TracNghiemLop11_Lan1
\begin{ex}%[1D3B2-2]%[ma:961a2f9d]
	Cho dãysố $\left(u_n\right)$ có $u_n=\dfrac{n+1}{2n+1}$. Số  $\dfrac{8}{15}$ là số hạng thứ bao nhiêu của dãy số $\left(u_n\right)$ ?
	\choice 
		{$8$} 
		{$6$} 
		 {$5$}
		 {\True $7$}
	 \loigiai{
	Giả sử $u_n=\dfrac{8}{15}\Rightarrow \dfrac{n+1}{2n+1}=\dfrac{8}{15}\Leftrightarrow 15n+15=16n=8\Leftrightarrow n=7$.} 
\end{ex}

% nguồn: Ok_TracNghiemLop11_Lan1
\begin{ex}%[1D3B2-2]%[ma:c0a61f1a]
	Cho dãy số $\left(y_n\right)$ xác định bởi $y_1=2$ và ${y}_{n+1}=2y_n+n^2-3n,\forall n\in \mathbb{N}^*$. Tổng $S_4$ của $4$  số hạng đầu tiên của dãy số là:
	 \choice
		 {\True $S_4=20$}
		 {$S_4=10$}
		{$S_4=30$} 
		{$S_4=14$} 
	\loigiai{ 
	 Ta tính được $y_2=2;y_3=4;y_4=12\Rightarrow S_4=20.$}
\end{ex}


%%%%%%%%%%%%%%%% ĐÚNG/SAI %%%%%%%%%%%%%%%%

% nguồn: Ok_dungsai_lop11_Lan1
\begin{ex}%[1D3N2-2]%[ma:6dbd39f3]
 Mỗi kết quả  giới hạn dưới đây đúng hay sai?
  \choiceTF
  {$\lim\limits_{x \to  1} (x-2)=1$}
  {$\lim\limits_{x \to  0} (x^2-x)=1$}
  {\True $\lim\limits_{x \to 3} \dfrac{x}{x+2}=\dfrac{3}{5}$} 
   {$\lim\limits_{x \to 2} \sqrt{x^2-1}=3$}
  \loigiai{
  \begin{itemchoice}
    \itemch $\lim\limits_{x \to 1} (x-2)=1-2=-1$.
   \itemch $\lim\limits_{x \to  0} (x^2-x)=1^2-1=0$.
   \itemch $\lim\limits_{x \to 3}  \dfrac{x}{x+2}=\dfrac{3}{3+2}=\dfrac{3}{5}$.
   \itemch  $\lim\limits_{x \to 2} \sqrt{x^2-1}=\sqrt{2^2-1}=\sqrt{3}$.
  \end{itemchoice}
 } 
\end{ex}

% nguồn: Ok_dungsai_lop11_Lan1
\begin{ex}%[1D3N2-2]%[ma:e1c1201a]
  Chọn phương án đúng hoặc sai cho các khẳng định dưới đây?
  \choiceTF
  {\True $\lim\limits_{x \to  0^+} \dfrac{1}{x}=+\infty$}
  {$\lim\limits_{x \to 1}\dfrac{x}{(x +  1)^2}=\dfrac{1}{2}$}
  {$\lim\limits_{x \to +\infty} \dfrac{2x +  1}{x^2 + 1}=1$}
   {\True $\lim\limits_{x \to 0} \dfrac{x}{\sqrt{x} + 1}=0$}
  \loigiai{
  \begin{itemchoice}
   \itemch Đúng\\ 
   Ta có $\lim\limits_{x\to 0^+} x=0$,  $\lim\limits_{x\to 0^+}1=1$ và $x>0$.\\
   Suy ra: $\lim\limits_{x \to 0^+}  \dfrac{1}{x}=+\infty$
    \itemch Sai\\
   Vì  $\lim\limits_{x \to 1}\dfrac{x}{(x + 1)^2}=\dfrac{1}{4}$.
   \itemch Sai\\ 
   Vì $\lim\limits_{x \to +\infty} \dfrac{2x + 1}{x^2 +  1}=0$.
    \itemch Đúng\\
   Vì  $\lim\limits_{x \to 0} \dfrac{x}{\sqrt{x} + 1}=\dfrac{0}{\sqrt{0}+1}=0$.
  \end{itemchoice}
 } 
\end{ex}

% nguồn: Ok_dungsai_lop11_Lan1
\begin{ex}%[1D3V2-2]%[ma:19241868]
 Các mệnh đề sau đúng hay sai? 
 \choiceTF 
  {Giới hạn $I_1=\lim \limits_{x  \rightarrow 0}(2 \sin x-1)=1$}
  {\True  Giới hạn $I_2=\lim \limits_{x \rightarrow \frac{\pi}{2}}(2 \cos 2 x+4)=2$}
  {Giới hạn $I_3=\lim \limits_{x \rightarrow-\frac{\pi}{3}}  2 \sin \left(x+\dfrac{2 \pi}{3}\right)=-\sqrt{3}$}
  {Giới hạn $I_4=\lim \limits_{x \rightarrow 0} \dfrac{\cos 2 x+(m+1) \cos x-m^2}{2  \sin ^2 x}=-\dfrac{7}{4}$. Khi đó $m \leq 0$}
  \loigiai{
  \begin{itemchoice}
   \itemch Sai. $I_1=\lim \limits_{x \rightarrow 0}(2 \sin x-1)=2 \sin 0-1=-1  $.
   \itemch Đúng. $I_2=\lim \limits_{x \rightarrow \frac{\pi}{2}}(2 \cos 2 x+4)=2 \cos \pi+4=2 .(-1)+4=2  $.
   \itemch Sai.$I_3=\lim \limits_{x \rightarrow -\frac{\pi}{3}} 2 \sin \left(x+\dfrac{2 \pi}{3}\right)=2 \sin \left(-\dfrac{\pi}{3}+\dfrac{2 \pi}{3}\right)=2 \sin \left(\dfrac{\pi}{3}\right)=2 \cdot \dfrac{\sqrt{3}}{2}=\sqrt{3} $. 
   \itemch Sai. Ta xét: $I_4=\lim \limits_{x \rightarrow 0} \dfrac{\cos 2 x+(m+1) \cos x-m^2}{2 \sin ^2  x}$.\\
   Do $\lim  \limits_{x \rightarrow 0}\left(2 \sin ^2 x\right)=0$ mà $I_4=-\dfrac{7}{4}$ nên giới hạn thuộc dạng vô định.\\
   Nên 
   \begin{align*}
    &\lim \limits_{x \rightarrow 0}\left(\cos 2  x+(m+1) \cos x-m^2\right)=0 \\
     \Leftrightarrow &1+(m+1) \cdot 1-m^2=0\\
     \Leftrightarrow&-m^2+m+2=0\\
    \Leftrightarrow &\hoac{& m=-1 \\ & m=2. }\\ 
   \end{align*}
   Khi $m=-1$, ta có $I_4=\lim \limits_{x \rightarrow 0} \dfrac{\cos  2 x-1}{2 \sin ^2 x}=\lim \limits_{x \rightarrow 0} \dfrac{-2 \sin ^2 x}{2 \sin ^2 x}=-1$.\\
   Khi $m=2$,  ta có
   \begin{align*}
    I_4=&\lim \limits_{x \rightarrow 0} \dfrac{\cos 2 x+3 \cos x-4}{2 \sin ^2  x}\\
    =&\lim \limits_{x \rightarrow 0} \dfrac{2 \cos ^2  x+3 \cos x-5}{2\left(1-\cos ^2 x\right)}\\
    =&\lim \limits_{x \rightarrow 0} \dfrac{(2 \cos x+5)(\cos x-1)}{-2(\cos  x+1)(\cos x-1)}\\
    =&\lim \limits_{x \rightarrow 0}  \dfrac{2 \cos x+5}{-2(\cos x+1)}\\
    =&\dfrac{2 \cdot 1+5}{-2  .(1+1)}\\
    =&-\dfrac{7}{4}. 
   \end{align*}
   Do đó, $m=2>0 $. 
  \end{itemchoice}
 } 
\end{ex}


%%%%%%%%%%%%%%%% TỰ LUẬN %%%%%%%%%%%%%%%%

% nguồn: Ok_TuanLuan_Toan11_Lan2
\begin{ex}%[1D3N2-2]%[ma:e3271bee]
	Giá trị của giới hạn $\lim\limits_{x\to  2}(3x^2+7x+11)$ là
	 \loigiai{
	$\lim\limits_{x\to 2}(3x^2+7x+11)=3\cdot 2^2+7\cdot 2+11=37$.  }
\end{ex}

% nguồn: Ok_TuanLuan_Toan11_Lan2
\begin{ex}%[1D3N2-2]%[ma:a5c9d856]
	Giá trị của giới hạn  $\lim\limits_{x\to 2}\dfrac{\sqrt[3]{3x^2-4}-\sqrt{3x-2}}{x+1}$ là
	\loigiai{ 
	Ta có $\lim\limits_{x\to 2}\dfrac{\sqrt[3]{3x^2-4}-\sqrt{3x-2}}{x+1}=\dfrac{\sqrt[3]{12-4}-\sqrt{6-2}}{3}=\dfrac{0}{3}=0$.  }
\end{ex}

% nguồn: Ok_TuanLuan_Toan11_Lan2
\begin{ex}%[1D3N2-2]%[ma:3dac93cc]
	Giá trị của giới  hạn $\lim\limits_{x\to 3}\sqrt{\dfrac{9x^2-x}{(2x-1)(x^4-3)}}$ là
	\loigiai{ 
	 $\lim\limits_{x\to 3}\sqrt{\dfrac{9x^2-x}{(2x-1)(x^4-3)}}=\sqrt{\dfrac{9\cdot 3^2-3}{(2\cdot 3-1)(3^4-3)}}=\dfrac{1}{\sqrt 5}$. }
\end{ex}

% nguồn: Ok_TuanLuan_Toan11_Lan2
\begin{ex}%[1D3N2-2]%[ma:e5e51156]
	Giá trị của giới hạn $\lim\limits_{x\to\sqrt 3}\left|x^2-4\right|$  là
	 \loigiai{
	$\lim\limits_{x\to\sqrt 3}\left|x^2-4\right|=\left|\left(\sqrt 3\right)^2-4\right|=1$.  }
\end{ex}

% nguồn: Ok_TuanLuan_Toan11_Lan2
\begin{ex}%[1D3N2-2]%[ma:43d69508]
	Giá trị của giới hạn $\lim\limits_{x\to 1}\dfrac{x-x^3}{\left(2x-1\right)\left(x^4-3\right)}$  là
	\loigiai{ 
	$\lim\limits_{x\to  1}\dfrac{x-x^3}{\left(2x-1\right)\left(x^4-3\right)}=\dfrac{1-1^3}{\left(2\cdot 1-1\right)\left(1^4-3\right)}=0$. }
\end{ex}

% nguồn: Ok_TuanLuan_Toan11_Lan2
\begin{ex}%[1D3N2-2]%[1D3H2-3]%[ma:c7bd9e4e]

	Tính các giới hạn sau đây 
	\begin{multicols}{2}
		\begin{enumerate}
			\item $\displaystyle\lim_{x \to  1} \left(x^2 + x - 2\right)$;
			\item $\displaystyle\lim_{x \to 0} \dfrac{x^2 + 3x  + 2}{x + 1}$;
			\item $\displaystyle\lim_{x \to -1} \dfrac{\sqrt{x + 4} -  1}{x + 3}$;
			\item $\displaystyle\lim_{x \to -2} \dfrac{x^2 + 5x + 2}{x^2  + x - 2}$;
			 \item $\displaystyle\lim_{x \to -1} \dfrac{x^2 + 3x + 2}{x + 1}$;
			\item $\displaystyle\lim_{x \to -3} \dfrac{x^2 - 9}{x  + 3}$.
		\end{enumerate}
	\end{multicols}
	\loigiai{ 
		\begin{enumerate}
			\item Ta có  $\displaystyle\lim_{x \to 1} \left(x^2+x-2\right)=1^2+1-2=0$.
			\item Ta có 
			\begin{eqnarray*}
				 \displaystyle\lim_{x \to 0} \dfrac{x^2+3x+2}{x+1}
				&=&  \dfrac{0^2+3 \cdot 0+2}{0+1}\\
				&=&\dfrac{2}{1} = 2. 
			\end{eqnarray*}
			\item Ta có 
			\begin{eqnarray*}
				\displaystyle\lim_{x \to -1}  \dfrac{\sqrt{x+4}-1}{x+3}
				&=& \displaystyle\lim_{x \to  -1} \left(\dfrac{\sqrt{x+4}-1}{x+3} \cdot \dfrac{\sqrt{x+4} + 1}{\sqrt{x+4}+1}\right) \\
				&=&  \displaystyle\lim_{x \to -1} \dfrac{x+3}{(x+3)\left(\sqrt{x+4}+1\right)} \\
				&=& \displaystyle\lim_{x  \to -1} \dfrac{1}{\sqrt{x+4}+1}\\
				&=&\dfrac{1}{\sqrt{3} + 1}. 
			\end{eqnarray*}
			\item  Ta có
			\begin{eqnarray*}
				\displaystyle\lim_{x \to -2} \dfrac{x^2+5x+2}{x^2+x-2} 
				 &=& \displaystyle\lim_{x \to -2} \dfrac{(x+2)(x+1)}{(x+2)(x-)} \\
				&=& \displaystyle\lim_{x  \to -2} \dfrac{x+1}{x-1}\\
				&=& \dfrac{-1}{-3}  = \dfrac{1}{3}.
			\end{eqnarray*}
			\item Ta có 
			\begin{eqnarray*}
				\displaystyle\lim_{x  \to -1} \dfrac{x^2+3x+2}{x+1}
				&=& \displaystyle\lim_{x  \to -1} \dfrac{(x+2)(x+1)}{x+1} \\
				 &=& \displaystyle\lim_{x \to -1} (x+2)\\
				&=&  1.
			\end{eqnarray*}
			\item  Ta có
			\begin{eqnarray*}
				\displaystyle\lim_{x  \to -3} \dfrac{x^2-9}{x+3}
				&=& \displaystyle\lim_{x \to  -3} \dfrac{(x-3)(x+3)}{x+3} \\
				&=&  \displaystyle\lim_{x \to -3} (x-3) \\
				&=& -6. 
			\end{eqnarray*}
		\end{enumerate}
	 }
\end{ex}

% nguồn: Ok_TuanLuan_Toan11_Lan2
\begin{ex}%[1D3N2-2]%[ma:005ed597]
	Giá trị của giới hạn  $\lim\limits_{x\to-1}\dfrac{\sqrt{3x^2+1}-x}{x-1}$ là
	 \loigiai{
	Ta  có $\lim\limits_{x\to-1}\dfrac{\sqrt{3x^2+1}-x}{x-1}=\dfrac{\sqrt{3+1}+1}{-1-1}=-\dfrac{3}{2}$. }
\end{ex}

% nguồn: Ok_TuanLuan_Toan11_Lan2
\begin{ex}%[1D3N2-2]%[ma:c5f58eb5]
	 Giá trị của giới hạn $\lim\limits_{x\to-1}\dfrac{x^2-3}{x^3+2}$ là
	\loigiai{ 
	 $\lim\limits_{x\to-1}\dfrac{x^2-3}{x^3+2}=\dfrac{(-1)^2-3}{(-1)^3+2}=-2$. }
\end{ex}

% nguồn: Ok_TuanLuan_Toan11_Lan1
\begin{bt}%[1D3Y2-2]%[ma:3e18c3a3]
	Cho dãy số $(u_n)$ có  số hạng tổng quát $u_n=\dfrac{2n+1}{n+2}$. Viết năm số hạng đầu của dãy số.
	\loigiai{ 
	Năm số hạng đầu của  dãy là\\ $u_1=1,u_2=\dfrac{5}{4},u_3=\dfrac{7}{5},u_4=\dfrac{3}{2},u_5=\dfrac{11}{7}$.}
\end{bt}

% nguồn: Ok_TuanLuan_Toan11_Lan1
\begin{bt}%[1D3Y2-2]%[ma:a6b54442]
	Cho dãy số $(u_n)$ được xác định bởi  $u_n=\dfrac{n^2+3n+7}{n+1}$. Viết năm số hạng đầu của dãy.
	 \loigiai{
	Ta có năm  số hạng đầu của dãy $u_1=\dfrac{1^2+3.1+7}{1+1}=\dfrac{11}{2}$, $u_2=\dfrac{17}{3},u_3=\dfrac{25}{4},u_4=7,u_5=\dfrac{47}{6}$.}
\end{bt}

% nguồn: Ok_TuanLuan_Toan11_Lan1
\begin{bt}%[1D3Y2-2]%[ma:aaeec134]
	Cho dãy số $(a_n)$ xác định bởi: $\heva{ 
	& u_1=-1,u_2=3  \\
	& {u}_{n+1}=5u_n-6{u}_{n-1} \forall n\ge 2 \\}$. Viết 7 số  hạng đầu tiên của dãy số.
	 \loigiai{
		 Bảy số hạng đầu của dãy là\\
		$u_3=5u_2-6u_1=21$; $u_4=5u_3-6u_2=87$;  $u_5=5u_4-6u_3=309$
	 $u_6=5u_5-6u_4=1023$; $u_7=5u_6-6u_5=3261$}
\end{bt}

% nguồn: Ok_TuanLuan_Toan11_Lan1
\begin{bt}%[1D3Y2-2]%[ma:4401ade4]
	 Cho dãy số $(u_n)$ xác định bởi $\heva{
	& u_1=1  \\
	& u_n=2{u}_{n-1}+3 \forall n\ge 2 \\}$. Viết năm số  hạng đầu của dãy số.
	 \loigiai{
		Ta có 5 số hạng  đầu của dãy là\\
	$u_1=1$;  $u_2=2u_1+3=5$; $u_3=2u_2+3=13$; $ u_4=2u_3+3=29$, $u_5=2u_4+3=61$.}
\end{bt}

% nguồn: Ok_TuanLuan_Toan11_Lan1
\begin{bt}%[1D3Y2-2]%[ma:00e29533]
	Cho dãy số  $(u_n):\heva{&u_1=u_2=1\\&u_n=u_{n-1}+u_{n-2}, \forall n\geq 3}$. Tìm số hạng thứ $7$ của dãy.
	 \loigiai{
		Ta có $u_3=u_1+u_2=2,  u_4=u_2+u_3=3, u_5=u_3+u_4=5, u_6=u_4+u_5=8, u_7=u_5+u_6=13$.
	 }
\end{bt}

% nguồn: Ok_TuanLuan_Toan11_Lan2
\begin{ex}%[1D3Y2-2]%[ma:b001bff3]
	Cho dãy số $(u_n)$ có số hạng tổng quát $u_n=\dfrac{2n+1}{n+2}$. Viết  năm số hạng đầu của dãy số.
	 \loigiai{
	 Năm số hạng đầu của dãy là\\ $u_1=1,u_2=\dfrac{5}{4},u_3=\dfrac{7}{5},u_4=\dfrac{3}{2},u_5=\dfrac{11}{7}$.}
\end{ex}

% nguồn: Ok_TuanLuan_Toan11_Lan2
\begin{ex}%[1D3Y2-2]%[ma:72224445]
	Cho  dãy số $(u_n)$ được xác định bởi $u_n=\dfrac{n^2+3n+7}{n+1}$. Viết năm số hạng đầu của dãy.
	 \loigiai{
	Ta có năm số hạng đầu của dãy $u_1=\dfrac{1^2+3.1+7}{1+1}=\dfrac{11}{2}$, $u_2=\dfrac{17}{3},u_3=\dfrac{25}{4},u_4=7,u_5=\dfrac{47}{6}$.} 
\end{ex}

% nguồn: Ok_TuanLuan_Toan11_Lan2
\begin{ex}%[1D3Y2-2]%[ma:55a810c0]
	Cho dãy số $(a_n)$  xác định bởi: $\heva{
	& u_1=-1,u_2=3 \\ 
	& {u}_{n+1}=5u_n-6{u}_{n-1} \forall n\ge 2 \\}$. Viết 7 số hạng đầu tiên của  dãy số.
	\loigiai{ 
		Bảy  số hạng đầu của dãy là\\
		 $u_3=5u_2-6u_1=21$; $u_4=5u_3-6u_2=87$; $u_5=5u_4-6u_3=309$
	$u_6=5u_5-6u_4=1023$; $u_7=5u_6-6u_5=3261$} 
\end{ex}

% nguồn: Ok_TuanLuan_Toan11_Lan2
\begin{ex}%[1D3Y2-2]%[ma:8ddc606a]
	Cho  dãy số $(u_n)$ xác định bởi $\heva{
	 & u_1=1 \\
	& u_n=2{u}_{n-1}+3 \forall n\ge 2 \\}$. Viết năm số hạng đầu của  dãy số.
	\loigiai{ 
		Ta có 5 số  hạng đầu của dãy là\\
	$u_1=1$; $u_2=2u_1+3=5$; $u_3=2u_2+3=13$; $  u_4=2u_3+3=29$, $u_5=2u_4+3=61$.}
\end{ex}

% nguồn: Ok_TuanLuan_Toan11_Lan2
\begin{ex}%[1D3Y2-2]%[ma:3967e915]
	Cho dãy số  $(u_n):\heva{&u_1=u_2=1\\&u_n=u_{n-1}+u_{n-2}, \forall n\geq 3}$. Tìm số hạng thứ $7$ của dãy.
	 \loigiai{
		Ta  có $u_3=u_1+u_2=2, u_4=u_2+u_3=3, u_5=u_3+u_4=5, u_6=u_4+u_5=8, u_7=u_5+u_6=13$.
	} 
\end{ex}
